% Generated draft: requires mathematical and source review.
\documentclass[11pt]{article}
\usepackage[margin=25mm]{geometry}
\usepackage{fontspec}
\setmainfont{Latin Modern Roman}
\usepackage{amsmath,amssymb,mathtools,unicode-math}
\setmathfont{Latin Modern Math}
\usepackage{enumitem,needspace,xurl,hyperref}
\setlength{\parindent}{0pt}
\setlength{\parskip}{5pt}
\setlength{\emergencystretch}{3em}
\newcommand{\N}{\mathbb{N}}
\newcommand{\Z}{\mathbb{Z}}
\newcommand{\Q}{\mathbb{Q}}
\newcommand{\R}{\mathbb{R}}
\newcommand{\C}{\mathbb{C}}
\newcommand{\F}{\mathbb{F}}
\newcommand{\A}{\mathbb{A}}
\newcommand{\PP}{\mathbb{P}}
\newcommand{\E}{\mathbb{E}}
\newcommand{\eps}{\varepsilon}
\newcommand{\dd}{\,\mathrm{d}}
\begin{document}
\section*{UnsolvedMath ID: 700049}
\textbf{Source problem code:} catalogue entry
\subsection{Definitions and mathematical statement}
For a finite positive integer \(r\), let \(F_r\) be the free group on \(r\) generators. For a group \(G\), write \(G'=[G,G]\) for its commutator subgroup and \(G''=[G',G']\). A group is metabelian if \(G''=\{1\}\). The free metabelian group of rank \(r\) is
\[
M_r=F_r/F_r''.
\]
A strict total order \(<\) on a group \(G\) is a transitive, irreflexive relation such that, for distinct \(g,h\in G\), exactly one of \(g<h\) and \(h<g\) holds. Such an order is invariant under multiplication on both sides if
\[
\forall a,b,g,h\in G,\qquad g<h\Longrightarrow agb<ahb.
\]
It is left invariant if \(g<h\Longrightarrow ag<ah\) for all \(a,g,h\in G\), and right invariant if \(g<h\Longrightarrow ga<ha\) for all \(a,g,h\in G\).

\textbf{Problem (A. I. Kokorin, 1965).} Describe all linear orderings of a free metabelian group with finitely many generators. Under the convention that a linearly ordered group has an order invariant under multiplication on both sides, this asks, for each finite positive rank \(r\), for a classification of
\[
\mathcal O_{\mathrm{bi}}(M_r)
=\{\,<\;:\;<\text{ is a strict total order on }M_r\text{ invariant under multiplication on both sides}\,\}.
\]
\textbf{Unresolved qualification.} The supplied record does not specify the invariance convention for ``linear orderings.'' If left-invariant or right-invariant orders are intended, the corresponding classification target uses that condition instead. The supplied excerpt does not resolve this distinction or specify whether orders are to be classified individually or up to an equivalence such as group automorphism. No rank restriction beyond finiteness is supplied; the positive ranks above cover the nontrivial free groups, and if rank zero is admitted, its free metabelian group is trivial and has a unique total order.
\subsection{Short English statement}
Kokorin asks for a description of all linear orderings of every finitely generated free metabelian group. The supplied statement leaves the precise order-invariance convention and any equivalence used in the classification unspecified.
\subsection{Sources}
\begin{itemize}
\item Local UnsolvedMath registry, permanent ID \(700049\), exact source problem code \texttt{catalogue entry}, title \emph{Kourovka Notebook Problem 1.54}; supplied complete source record.
\item \emph{The Kourovka Notebook (Unsolved Problems in Group Theory)}, original Problem 1.54, first issue (1965); proposed by A. I. Kokorin. The supplied background identifies the reference edition as the 21st issue edition (2026), \texttt{arXiv:1401.0300v48}.
\item The supplied collection description attributes the edition covering Issues 1--21 (1965--2026) to editors E. I. Khukhro and V. D. Mazurov. The local record labels this problem \emph{open}; that label is supplied curation and has not been independently verified here.
\end{itemize}
\end{document}
